\section[Application of the purity theorem]{Application of the purity theorem:
continuity theorem for the fundamental groups of the
fibers of a proper and smooth morphism}
\label{X.3}

Let us recall without proof\footnote{For a proof, \cf SGA~2 X~\Ref{X.3.4}.} the
following

\begin{puritytheorem}[Zariski--Nagata]
\label{X.3.1}
Let $f\colon X\to Y$ be a
quasi-finite and dominant morphism of integral preschemes, with $X$
normal, $Y$ regular locally noetherian, and let $Z$ be
the set of points of~$X$ where $f$ is not \'etale,
\ie where~$f$ is ramified (this amounts to the same thing,
\textup{I~\Ref{I.9.5}~(ii)}). If $Z\ne X$, $Z$ is of codimension~$1$ in
$X$ at all its points, \ie for every irreducible component
$Z'$ of~$Z$ with generic point~$z$, the Krull dimension of
$\cal{O}_{X,z}$ is equal to~$1$.
\end{puritytheorem}

Recall that a prescheme is said to be \emph{normal}
\resp \emph{regular} if its local rings are normal
\resp regular, and that the relation $Z\ne X$ also means that
the finite extension $R(Z)/R(X)$ (where $R$ denotes the field of
rational functions) is \emph{separable}. Placing oneself at
the generic point $z$ of a component $Z'$ of~$Z$, and
localizing at the point $y$ of~$Y$ below~$z$, one finds
% original p. 276
the equivalent statement:

\begin{corollary}
\label{X.3.2}
Let $A$ be a regular noetherian local ring, $A\to B$ an
injective local homomorphism such that $B$ is normal, a localization
of an algebra of finite type over~$A$, and \emph{quasi-finite} over~$A$;
suppose moreover that $\dim A (=\dim B)\ge 2$, and that for every
prime ideal $\goth{p}$ of~$B$ distinct from the maximal ideal,
$B$ is \'etale over~$A$ at~$\goth{p}$, \ie $B_\goth{p}$ is
\'etale over~$A_\goth{q}$ (where $\goth{q}=A\cap\goth{p}$). Then
$B$ is \'etale over~$A$.
\end{corollary}

Moreover, it is not difficult to reduce this last
statement to the case where $A$ is a \emph{complete} local ring,
hence where $B$ is \emph{finite} over~$A$. Zariski~\cite{X.5} gives a
simple proof of this result, valid in the case
of equal characteristics; the general case is due
to Nagata~\cite{X.3}, who relies on a delicate result of
Chow; the latter has not been checked by any of the
participants of the Seminar, and should be the subject of a
later expos\'e. Let us only point out here the
very simple proof in the special case where $\dim A=2$, which is
sufficient for the most important application that we shall make of it
in the present no. Since $B$ is normal, it is a
$B$-module of depth (old terminology: cohomological
codimension) $\ge 2$, hence it is an $A$-module of depth~$\ge
2$, and since $A$ is regular of dimension $2$, it follows
that $B$ is a \emph{free module} over~$A$\kern1pt\footnote{\Cf EGA $0_{\textup{IV}}$
17.3.4}. It then follows from I~\Ref{I.4.10} that the set
of prime ideals~$\goth{q}$ of~$A$ at which $B$ is
ramified over~$A$ is the subset of~$\Spec(A)$ defined by a
principal ideal (generated by the discriminant of a basis of
$B$ over~$A$), hence is empty if it is contained in the closed
point of~$\Spec(A)$, which proves~\Ref{X.3.2} when $\dim A=2$.

We shall use~\Ref{X.3.1} above all in the equivalent form:

\begin{corollary}
\label{X.3.3}
Let $X$ be a regular locally noetherian
prescheme, $U$ an open subset of~$X$ complementary
to a closed subset $Z$ of~$X$ of codimension~$\ge 2$. Then the
functor $X'\mto X'\times_X U$ from the category of
\'etale coverings of~$X$ to the category of
\'etale coverings of~$U$ is an equivalence
% original p. 277
of categories; in particular, if $a$ is a
geometric point of~$U$, the canonical homomorphism $\pi_1(U,a)\to
\pi_1(X,a)$ is an isomorphism.
\end{corollary}

The last assertion is evidently a consequence of the
first, and for the latter one may evidently assume that $X$
is connected, hence irreducible. From the normality of~$X$
it already follows that the functor $X'\mto X'\times_X U$ from
the category of locally free coverings (not
necessarily \'etale) of~$X'$ to the category of
coverings of~$U$ is fully faithful, since the functor
$\cal{E}\mto\cal{E}\vert U$ from the category of locally free
Modules on~$X$ to the category of locally
free Modules on~$U$ is so. It therefore remains to prove that for every
\emph{\'etale} covering $U'$ of~$U$, there exists an
\'etale covering $X'$ of~$X$ (necessarily unique
by what precedes), such that $U'$ is isomorphic to
$X'\times_X U$. One may evidently assume $U'$ connected, hence
irreducible since ($U$~being normal) $U'$ is normal. Let
$K$ be the field of rational functions on~$X$, or on~$U$ (it is
the same), $K'$ that of~$U'$; then $U'$ is identified with the normalization
of~$U$ in~$K'$ I~\Ref{I.10.3}. Let $X'$ be the normalization of
$X$ in~$K'$ (EGA~II~6.3); then $X'=X\times_X U \cong U'$; on the other
hand $X'$ is normal, integral, and the structural morphism $f\colon
X'\to X$ is \emph{finite} and dominant (since $X$ is normal and $K'/K$ is
a finite separable extension). It is \'etale in
$U'=f^{-1}(U)=X'=f^{-1}(Z)$, and since $Z$ is of codimension $\ge 2$
in $X$, $f^{-1}(Z)$ is of codimension $\ge 2$ in~$X'$. One then concludes
from~\Ref{X.3.1} that $X'$ is \'etale over~$X$, which completes
the proof.

\medskip
Now let $f\colon X\to Y$ be a rational map from a
locally noetherian and regular prescheme $X$
to a prescheme~$Y$, and suppose that $f$ is defined
in an open set $U$ complementary to a closed subset of
codimension~$\ge 2$. Then one deduces from~\Ref{X.3.3} a functor,
defined up to isomorphism, from the category of
\'etale coverings of~$Y$ to the category of
\'etale coverings of~$X$, whence for every
geometric point $a$ of~$U$, with image $b$ in~$Y$, a canonical
homomorphism
$$
\pi_1(X,a)\to \pi_1(Y,b)
$$
(deduced
% original p. 278
from the canonical homomorphism $\pi_1(U,a)\to \pi_1(Y,b)$ by means
of the isomorphism $\pi_1(U,a)\isomto \pi_1(X,a))$. When $f$ is
a dominant morphism, $X$ and $Y$ being integral with fields $K$
and~$L$, so that $K$ is an extension of~$L$, and when $Y$ is
normal, these correspondences can be made precise in terms of field
extensions by noting that for every finite extension $L'$ of~$L$,
unramified over~$Y$, the algebra $K'=L'\otimes_L K$ over~$K$ is
unramified over~$X$.

In particular, these reflections show that the fundamental group
of connected regular locally noetherian preschemes,
pointed by geometric points
localized in codimension~$\le 1$, is a \emph{functor} when
one takes as morphisms in this category the dominant rational
maps defined in complements of
closed subsets of codimension~$\ge 2$. Recalling for example
that a rational map from a normal scheme over a field $k$
to a scheme proper over~$k$ is defined in the
complement of a set of codimension~$\ge 2$, one finds:

\begin{corollary}[Birational invariance of the fundamental group]
\label{X.3.4}
Let $k$ be a field, $X$ and $Y$ two schemes proper over~$k$ and
regular, $f\colon X\to Y$ a birational map from~$X$
to~$Y$, $\Omega$ an algebraically closed extension of the
function field $K$ of~$X$, allowing one to define the fundamental group
of~$X$ and the fundamental group of~$Y$. The latter are then
canonically isomorphic.
\end{corollary}

This also means that for a finite extension $K'$ of~$K$, if it
is unramified over a proper nonsingular ``model'' $X$
of~$K$, it is so over every other proper nonsingular model.

\begin{remark}
\label{X.3.5}
For other applications of the purity theorem, see the
work
of \textsc{Abhyankar}
expounded in~\cite{X.4}, inspired by the
results of \textsc{Zariski} \cite[chap\ptbl VIII]{X.6}, proved by
topological means. The latter are far from having been assimilated
by ``abstract'' algebraic geometry and deserve
new efforts.
\end{remark}

We
% original p. 279
shall need a few elementary facts from the theory
of ramification. Let $V$ be a discrete valuation ring with
field of fractions~$K$, residue field~$k$, $L$ a Galois
extension of~$K$ with group~$G$, $V'$ the normalization of~$V$
in~$L$, which is a free module of rank $n=[L{\colon }K]$ over~$V$,
$\goth{m'}$ a maximal ideal of~$V'$, $G_d$ the subgroup of~$G$
formed by the elements leaving $\goth{m'}$ invariant, so
that $G_d$ operates on the residue extension
$k'=V'/\goth{m'}$ of~$k$, and $G_i$ the subgroup of the elements
of~$G_d$ operating trivially (recall that $G_d$ and $G_i$ are
called respectively the \emph{decomposition} and
\emph{inertia} subgroups of~$G$). One says that $L$ is \emph{tamely
ramified} over~$V$ if $n_i=[G_i{\colon }e]$ is of order $p$ prime
to the characteristic of~$k$ (a condition always
satisfied if $k$ is of characteristic~0). It is well
known that $G_i$ then embeds canonically in the group ${k'}^*$,
hence is isomorphic to the group of
$n_i$th roots of unity
in~${k'}^*$, which implies in particular that $G_i$ \emph{is
cyclic}. The typical case of this situation is the one where one sets
$L=K[t]/(t^n - u)$, $u$ being a uniformizer of~$V$ and $n$ an
integer prime to~$p$: if $K$ contains the $n$th roots of
unity, $L$ is a totally ramified Galois extension
of~$K$, with Galois group $G=G_i$ isomorphic to~$\ZZ/n\ZZ$.

\begin{lemma}[\textsc{Abhyankar}'s lemma]
\label{X.3.6}
Let $V$ be a discrete valuation ring with
field of fractions~$K$, $L$ and $K'$ two Galois extensions
of~$K$ \emph{tamely ramified} over~$V$, $n$ and $m$ the
orders of the
corresponding inertia groups,
$L'$ a composite extension
of~$L$ and $K'$ over~$K$. If $m$ is a multiple of~$n$,
then $L'$ is unramified over the localizations of the
normal closure $V'$ of~$V$ in~$K'$.
\end{lemma}

Indeed, let $W'$ be the normalization of~$V'$ in~$L'$, $\goth{m'}$
a maximal ideal of~$V'$, $\goth{n'}$ a maximal ideal of~$W'$
above
$\goth{m'}$,
$\goth{n}$ the maximal ideal it induces on the normalization $W$
of~$V$ in~$L$, $G,\ H,\ M$ the Galois groups of~$L,\ K',\ L'$
over~$K$, and $G_i,\ H_i,\ L_i'$ the inertia groups corresponding to the
chosen maximal ideals. Then $M$ embeds in the product
$G\times H$ and $M_i$ in the product $G_i\times H_i$, in such a way that the projections $M\to G$ and $M\to H$, $M_i\to G_i$
% original p. 280
and $M_i\to H_i$ are surjective (sorites of the intermediate
field). It already follows, since $G_i$ and
$H_i$ are by hypothesis cyclic of orders $m$ and $n$ prime
to~$p$, that $M_i$ is of order prime to~$p$, hence cyclic, and
since $m$ is a multiple of~$n$, so that the elements of~$G_i\times
H_i$ have zero
$m$th power, $M_i$ is of order
dividing~$m$, hence of order equal to $m$ since $M_i\to H_i$ is
surjective. This last homomorphism is therefore also
injective. Now its kernel is the inertia group of~$\goth{n'}$
above~$\goth{m'}$, which proves that $L'$ is unramified
over~$K'$ at~$\goth{n'}$. Whence the lemma.

Let us now place ourselves under the conditions of~\Ref{X.2.4},
where one has a specialization homomorphism
$$
\pi_1(\overline X_1,\overline a_1)\to\pi_1(\overline X_0,\overline
a_0)
$$
which is \emph{surjective}, relative to a proper and
separable morphism $f\colon X\to Y$. We wish to determine more precisely the kernel of
this homomorphism. Proceeding as in the proof
of~\Ref{X.2.4}, one sees that in this question one may always
assume that $Y$ is the spectrum of a \emph{discrete valuation
ring~$V$, complete and with algebraically closed
residue field} (for one can always find such a ring
and a morphism of its spectrum $Y'$ into $Y$ whose image is
$\{y_0,y_1\}$). Then one has $X_0=\overline X_0$, $\kres(y_0)=k=$
residue field of~$V$, $\kres(y_1)=K=$ field of fractions of~$V$. Let
$K_s$ be the separable closure of~$K$, $\overline K$ its
algebraic closure, and for every subring $W$ of~$\overline
K$ containing~$V$, set $X_W=X\otimes_V W$. In particular one has
$$
X_V=X, \qquad X_K=X_1, \qquad X_{\overline K}=\overline X_1.
$$
Moreover the canonical morphism $\overline X_1=X_{\overline K}\to
X_{K_s}$ induces an isomorphism on the fundamental groups
(IX~\Ref{IX.4.11}), so that, taking into account
the isomorphism~\Ref{X.2.1} $\pi_1(X_0)\to\pi_1(X)$, one is reduced
to studying the surjective homomorphism
$$
\pi_1(X_{K_s})\to\pi_1(X)
$$
associated
% original p. 281
with the canonical morphism $X_{K_s}\to X$. The determination of the kernel
of the latter amounts to the solution of the following problem:
\emph{one has a connected principal covering $Z_{K_s}$ of
$X_{K_s}$, with group~$G$} (hence associated with a homomorphism of
$\pi_1(X_{K_s})$ into~$G$), \emph{determine under what
conditions it is isomorphic to the inverse image of a
principal covering $Z$ of~$X$ with group~$G$.}

Let us note first that $K_s$ is the filtered increasing union of its
finite subextensions $K'$ over~$K$, and that consequently $Z_{K_s}$
is isomorphic to the inverse image of a principal covering
$Z_{K'}$ of~$X_{K'}$ for a suitable $K'$ (one should be careful,
however, that for fixed $K'$, $Z_{K'}$ is not uniquely
determined). To say that $Z_{K_s}$ is isomorphic to the inverse
image of a principal covering $Z$ of~$X$ means that there exists
a finite subextension $K''\supset K'$ of~$K_s$ such that
$Z_{K''}=Z_{K'}\otimes_{K'} K''$ is isomorphic to $Z\otimes_V
K''$. Let us now denote, for a finite subextension $K'$
of~$K_s$, by $V'$ the normalization of~$V$ in~$K'$, which is a
complete discrete valuation ring with
residue field~$k$. Hence the canonical morphism $X_{V'}\to X_V$ induces
an isomorphism for the fibers above the closed points of
$Y=\Spec(V)$ and $Y'=\Spec(V')$, and it then follows from~\Ref{X.2.1}
applied to $X_V$ and $X_{V'}$ that the induced homomorphism on
the fundamental groups $\pi_1(X_{V'})\to \pi_1(X_V)$ is an
isomorphism, or again that every principal covering of~$X_{V'}$
is the inverse image of a principal covering of~$X_V$
determined up to isomorphism. This therefore implies the
following

\begin{lemma}
\label{X.3.7}
Let $Z_{K'}$ be a connected principal covering of~$X_{K'}$ with
group~$G$, $Z_{K}$ its inverse image on~$X_{K_s}$\quoi. For the
latter to be isomorphic to the inverse image of a
principal covering $Z$ of~$X$, it is necessary and sufficient that there exist a
finite extension $K''\supset K'$ of~$K$ in $K_s$ such that the
principal covering $Z_{K''}$ of~$X_{K''}$ is induced by a
principal covering of~$X_{V''}$.
\end{lemma}

Suppose in particular that the $X_{V''}$ are normal (for this it
suffices for example that $X_0$ be normal, and a fortiori that $X_0$
be simple, \cf I~\Ref{I.9.1}).
% original p. 282
Since they are connected, they are therefore irreducible. Let $L$ be the
field of rational functions for $X$ and $X_{K}$, $L'$ the one for
$X_{V'}$ and $X_{K'}$, $L''$ the one for $X_{V''}$ and~$X_{K''}$\quoi. Then
under the conditions of~\Ref{X.3.7}, $Z_{K'}$ defines a finite
separable extension $R'$ of~$L'$, and $Z_{K''}$ defines the extension
$R''=R'\otimes_{L'} L'' = R'\otimes_{K'} K''$. The condition
considered in~\Ref{X.3.7} therefore also means that there exists a
finite separable extension $K''$ of~$K'$ such that
$R''=R'\otimes_{K'} K''$ is \emph{unramified} over the
normal scheme $X_{V''}$ with field $L''=L'\otimes_{K'} K''$, and not
only over the open subset $X_{K''}$ of~$X_{V''}$.

We assume from now on that $f\colon X\to Y$ is a
\emph{smooth} morphism; hence the morphisms $X_{V'}\to \Spec(V')$ are smooth,
hence the schemes $X_{V'}$ are \emph{regular.} Note that the
fiber of the closed point of~$\Spec(V')$ in $X_{V'}$ is
irreducible and of codimension~$1$. Let $\goth{o'}$ be its local
ring, which is therefore a discrete valuation ring with field
$L'$, with residue field isomorphic to the field of rational
functions of~$X_0$, hence having the same characteristic as
$k$. Let us define similarly $\goth{o''}$ in $L''$, which is
evidently the normalization of
$\goth{o'}$
in~$L''$. It
then follows from the purity theorem~\Ref{X.3.1}
or~\Ref{X.3.3} that for $R''$ to be unramified over~$X_{V''}$,
it is necessary and sufficient that $R''$ be unramified over~$\goth{o''}$,
the normalization of~$\goth{o'}$ in~$L''$.

Let us now note that if $u'$ is a uniformizer of~$V'$, it is
also a uniformizer of~$\goth{o'}$. If then $n$ is an integer
prime to the characteristic $p$ of~$k$, and if one takes
$K''=K'[t]/(t^n-u)$, then
$K''$
is a finite Galois extension
of~$K'$ and $L''$ is isomorphic to $L'[t]/(t^n-u')$, hence is
tamely ramified over~$\goth{o'}$ with inertia group
of order~$n$. Suppose then $G$ \emph{of order prime to~$p$,} which
implies that $R'$ is tamely ramified over~$\goth{o'}$, and let us take for $n$ a multiple prime to $p$ of
the order of the inertia group of~$R'$ over~$\goth{o'}$ (for example
$n=[G\colon e]$). Applying
\textsc{Abhyankar}'s lemma~\Ref{X.3.6}, one
sees that the condition considered in~\Ref{X.3.7} is
satisfied.

This proves the following theorem:

\begin{theorem}
\label{X.3.8}
Let $f\colon X\to Y$ be a proper and smooth morphism, with
% original p. 283
geometrically connected fibers, with $Y$ locally noetherian,
$y_0$ and $y_1$ two points of~$Y$ such that
$y_0\in\overline{\{y_1\}}$,
$\overline{X}_0$ and $\overline{X}_1$ the corresponding geometric
fibers; consider the specialization homomorphism
\eqref{X.2.4} $\pi_1(\overline X_1)\to\pi_1(\overline X_0)$. This
homomorphism is surjective, and every continuous homomorphism of
$\pi_1(\overline X_1)$ into a finite group $G$ of order prime to
the characteristic $p$ of~$\kres(y_0)$ comes from a homomorphism of
$\pi_1(\overline{X}_0)$ into~$G$.
\end{theorem}

In other words:

\begin{corollary}
\label{X.3.9}
If $\kres(y_0)$ is of characteristic zero, then the specialization
homomorphism is an isomorphism. If
$\kres(y_0)$ is of characteristic $p>0$,
then the kernel of
the specialization homomorphism is contained in the intersection
of the kernels of the continuous homomorphisms of~$\pi_1(\overline X_1)$ into
finite groups of order prime to~$p$ (or again, the closed invariant
subgroup generated by a Sylow $p$-subgroup of the
group of Galois type $\pi_1(\overline X_1)$); if therefore
$\pi_1(\overline X_1)^{(p)}$
denotes the quotient group of
$\pi_1(\overline X_1)$ by the preceding closed subgroup,
and if one defines $\pi_1(\overline X_0)^{(p)}$ similarly, then
the specialization homomorphism induces an \emph{isomorphism}
$$
\pi_1(\overline X_1)^{(p)}\isomto \pi_1(\overline X_0)^{(p)}
$$
\end{corollary}

One will note that the proof of \Ref{X.3.8} is purely algebraic.
Proceeding as in \Ref{X.2.6}, one concludes from it \emph{by
transcendental means:}

\begin{corollary}
\label{X.3.10}
Let $X_0$ be a proper, smooth and connected curve of genus $g$ over an
algebraically closed field of characteristic~$p$. With
the notation introduced in~\Ref{X.3.9}, the group $\pi_1(X_0)^{(p)}$
is isomorphic to $\Gamma^{(p)}$, where $\Gamma$ is the group of
Galois type generated by generators $s_i,~t_i$
($1\le i\le g$) subject to the relation
$$
(s_1t_1s_1^{-1}t_1^{-1})\cdots (s_gt_gs_g^{-1}t_g^{-1})=1
$$
\end{corollary}

\begin{remarks}
\label{X.3.11}
% original p. 284
In the case where $\kres(y_0)$ is of characteristic zero, the
result~\Ref{X.3.9} is well known by transcendental means. One
will note that the proof of~\Ref{X.3.10} makes use of the
purity theorem in the case of unequal
characteristics, but only in the case of rings of dimension 2,
where the proof of the said theorem is
easy and was recalled in the text.
\end{remarks}

\begin{thebibliography}{0}{X.4}

\bibitem[1]{X.1} K\ptbl \textsc{Kodaira}, On compact analytic surfaces,
Ann.\ of Math.\ \textbf{71} (1960), p\ptbl 111--152.

\bibitem[2]{X.2} S\ptbl \textsc{Lang} and
J-P\ptbl\textsc{Serre}, Sur les rev\^etements non
ramifi\'es des vari\'et\'es alg\'ebriques,
Amer.\ J.\ Math.\ \textbf{79} (1957), p\ptbl319--330.

\bibitem[3]{X.3} M\ptbl \textsc{Nagata},
On the purity of branch loci in regular local rings,
Illinois J.\ Math. \textbf{3} (1959), p\ptbl 328--333.

\bibitem[4]{X.4}
J-P\ptbl\textsc{Serre}, Rev\^etements ramifi\'es du plan
projectif (d'apr\`es S\ptbl Abhyankar), S\'eminaire Bourbaki, May
1960.

\bibitem[5]{X.5} O\ptbl \textsc{Zariski},
On the purity of the branch locus of algebraic functions,
Proc.\ Nat.\ Acad.\ Sc.\ USA  \textbf{44} (1958), p\ptbl 791--796.

\bibitem[6]{X.6} O\ptbl \textsc{Zariski},
Algebraic surfaces, Ergebnisse... 1948, Chelsea (New York).
\end{thebibliography}
